% Authoritative LaTeX manuscript; Markdown is a generated compatibility projection.

% <!-- part-title: The Theory -->

% <!-- chapter-title: Testing the Theory -->

\chapter{Testing the Theory}

% <!-- point: comparisons-test-propositions-and-guide-investment | Shared measurements and experimental designs support tests of all four propositions and decisions about engineering investment. -->

The theory developed here began with one longitudinal study, supplemented by the industrial accounts examined in \magevalue{chapter:software-factories}. The preceding sections have used that evidence to develop four propositions about representation, enforcement, engineering capital, and agentic capacity. Testing those propositions requires comparisons across tasks, systems, engineering environments, and agents, including observations that might contradict the predicted relationships.

The comparisons also matter to engineering practice. The organizations examined in \magevalue{chapter:software-factories} differ in the knowledge they preserve, the controls they enforce, and the capabilities available to them. A useful representation in one environment may duplicate knowledge already available in another; a new validator may reduce costly failures or merely add friction. Measuring these effects can guide which parts of MAGE an organization adopts, where it invests next, and when an existing arrangement should be reconsidered. This section develops the measurements, experimental designs, and engineering decisions that follow from the four propositions.

\section{Measuring Engineering Progress}\label{heading:research-measuring-engineering-progress}

Lines of code, commits, and token expenditure cannot by themselves establish engineering progress. Agents can produce large quantities of code with little human effort, while substantial engineering work may produce no code. Tokens measure resources consumed: Agarwal and colleagues found wide variation in token expenditure across weeks with similar commit counts, and engineers reported using agents for analysis, research, and planning as well as implementation~\citep{agarwal2026tokens}.

\autoref{engineering-measures} distinguishes candidate measures of implementation, engineering models, the GEE, and engineering activity. Model-based engineering offers relevant precedents in its measures of model maturity, traceability, verification, rework, and engineering effort~\citep{garcia2021modelproductivity,henderson2023digitalmetrics,campo2023mbse}. These categories describe different aspects of output, capability, and activity; they are not components of a prescribed productivity score.

{\footnotesize
\begin{longtable}[]{@{}
  >{\raggedright\arraybackslash}p{(\linewidth - 4\tabcolsep) * \real{0.20}}
  >{\raggedright\arraybackslash}p{(\linewidth - 4\tabcolsep) * \real{0.54}}
  >{\raggedright\arraybackslash}p{(\linewidth - 4\tabcolsep) * \real{0.26}}@{}}
\caption[Candidate measures of agentic engineering]{\emph{Candidate measures of agentic engineering.} Implementation measures describe changes to the product; modeling measures describe explicit engineering knowledge; GEE measures describe the machinery supporting controlled change; activity measures describe resources consumed and work performed. Examples are illustrative rather than exhaustive.}\label{engineering-measures}\tabularnewline
\toprule\noalign{}
Measurement target & Candidate measures & What they characterize \\
\midrule\noalign{}
\endfirsthead
\toprule\noalign{}
Measurement target & Candidate measures & What they characterize \\
\midrule\noalign{}
\endhead
\bottomrule\noalign{}
\endlastfoot
Implementation & LOC added, deleted, or retained; code churn; commits and merged changes; completed features; escaped defects; rework; accepted changes per unit time & Product change, delivery, and durability \\
Engineering models & Model elements or size; model maturity; represented requirements and properties; model coverage; relationships across views or abstraction levels; traceability completeness; properties analyzed or validated & Explicit engineering knowledge, relationships, and evidence \\
Governed engineering environment (GEE) & Number and coverage of constraints, sensors, and gates; obligations with independent checks; model–implementation correspondence; violations detected or blocked; false rejections; human interventions; control maintenance cost; changes admitted without human inspection & Capacity to direct, evaluate, and govern implementation \\
Engineering activity and resources & Agent invocations; tokens consumed; inference expenditure; agent execution time; human engineering and review time; dispatches; retries; tool calls; validation runs & Work performed and resources consumed \\
\end{longtable}
}

The categories need not move together. A migration may change thousands of lines while preserving externally visible behavior. A model can make an architectural obligation explicit without changing the implementation; a gate can prevent violations without adding a feature. Conversely, an agent may complete substantial implementation without improving the environment. The engineering question determines which measures matter: accepted changes for implementation capacity, represented properties and analyses for Modeling, or coverage and intervention for Alignment.

Durability also has more than one object. Product functionality may survive subsequent changes, while a model continues to answer consequential questions or a control remains effective as the system evolves. These are different forms of durable engineering progress. Their size alone does not establish their value; subsequent use and the costs of sustaining that use matter.

\autoref{theory-experimental-comparisons} relates these measures to the four propositions. Each comparison holds the proposed intervention against an alternative that could plausibly work.

{\footnotesize
\begin{longtable}[]{@{}
  >{\raggedright\arraybackslash}p{(\linewidth - 4\tabcolsep) * \real{0.22}}
  >{\raggedright\arraybackslash}p{(\linewidth - 4\tabcolsep) * \real{0.43}}
  >{\raggedright\arraybackslash}p{(\linewidth - 4\tabcolsep) * \real{0.35}}@{}}
\caption[Comparisons that distinguish MAGE's proposed effects]{Four comparisons connecting the theoretical propositions to observable outcomes.}\label{theory-experimental-comparisons}\tabularnewline
\toprule\noalign{}
Proposition & Comparison & What to measure \\
\midrule\noalign{}
\endfirsthead
\toprule\noalign{}
Proposition & Comparison & What to measure \\
\midrule\noalign{}
\endhead
\bottomrule\noalign{}
\endlastfoot
Representation leverage (P1) & Structured representation versus prose containing the same obligations; separate elicitation comparison & Correctness, obligation completeness, and reasoning cost \\
Representation-enabled enforcement (P2) & Model-derived enforcement versus producer guidance alone, with existing checks held fixed & Escape, false rejection, rejection cost, and correct delivery \\
Governance conversion and engineering capital (P3) & Local repair or on-demand reconstruction versus reusable models, procedures, or controls & Later capability, recurrence, and full cost, including construction and upkeep \\
Governable acceleration (P4) & Agentic capacity crossed with specified representations, controls, or evidence pipelines & Durable throughput, defect escape, and human attention \\
\end{longtable}
}

Independent assessment of accepted work and the lifecycle accounting in \magevalue{sec:toward-a-theory-of-mage} provide a common basis. Comparisons must include failed attempts, repair, human attention, and upkeep, report unresolved requests, and use declared quality requirements and evaluation horizons. This permits an early construction cost to be compared with later savings without concealing work that was never delivered.

\section{Designing the Comparisons}\label{heading:research-designing-the-comparisons}

Controlled comparisons hold tasks, acceptance criteria, and relevant resources fixed while varying an intervention. Random assignment and repeated trials, where feasible, help distinguish its effect from differences among tasks and agent executions. Factorial comparisons vary factors independently to reveal their interaction: a representation may substitute for a stronger reasoner, or become more useful when a stronger reasoner can exploit it. The mechanism under test must not be the sole judge of whether its results are acceptable. All conditions need independent outcome assessment.

Longitudinal comparisons follow the introduction, reuse, maintenance, and retirement of engineering assets. Staged introductions across comparable subsystems or teams can help distinguish an intervention from changes already occurring over time. Such studies are needed when an investment changes later work rather than merely the next task. Teams, task families, or repositories within one organization can supply comparisons, provided differences in difficulty, personnel, obligations, and existing engineering knowledge are accounted for.

GEE conditions should differ in specified representations, controls, or evidence pipelines, with their adequacy assessed independently of the measured outcomes. The dimensions established in \autoref{6.1-toward-a-theory-of-mage:heading:the-dynamic-model} guide that characterization. A condition cannot be called better simply because it produces better results. Each study can begin with a particular engineering problem, such as repeated reconstruction or expensive review, and use the relevant proposition to identify the relationship to test.

\section{Testing Representation Leverage (P1)}\label{heading:representation-engineering}

\emph{Hypothesis H1 (Representation leverage). For tasks requiring substantial reconstruction of system knowledge, a trustworthy structured representation reduces the cost of obtaining an independently acceptable result compared with equivalent information supplied in prose.}

% <!-- point: natural-language-descriptions-are-models-so-compare-against-strong-baselines | Compare structured representations with equivalent prose to isolate form; study elicitation separately to learn whether engineers encode better obligations. -->

One experiment holds the obligations, task, foundation model, tools, and admission mechanism fixed while varying representation form. Measure correct delivery, attempts, reconstruction effort, inference cost, and human interpretation or repair. Include construction and maintenance over the declared horizon. This tests fulfillment under the same encoded obligations. If structured models consistently add cost without improving accepted delivery, H1 would be weakened for those task families.

A separate elicitation experiment should test whether structured representations help engineers discover or express obligations they would otherwise omit. Engineers could receive the same problem and background evidence but use different representational forms; independent domain assessment would examine the resulting obligations before realization. This tests the encoding contribution without assuming that treatments begin with equally complete specifications. The two experiments distinguish a useful form from additional information supplied only to one condition.

% <!-- point: experiment-compares-authored-induced-and-innovated-representations | Compare designed, recovered and invented representations using independent correspondence evidence and held-out engineering questions, including costs and later benefit. -->

Further experiments can compare engineer-designed representations, representations recovered from implementation, and agent-proposed abstractions, using raw implementation as a common baseline. Recovery should expose consequential structure already present; innovation should make useful relationships available that existing representations do not. Held-out engineering questions and failures can test usefulness beyond the demonstration. Assess claimed correspondence independently, along with human comprehensibility, reconstruction cost, reasoning success, and maintenance. The comparison should also ask whether the structure is more economical to retain or reconstruct when needed.

% <!-- point: representation-has-oversight-value-beyond-realization | Representation studies should test effects on human engineering decisions as well as production; the review experiment under P4 examines one such use. -->

Representations can also change human engineering decisions. The review experiment under P4 examines whether models and evidence help engineers judge the growing volume of agent-produced work.

\section{Testing Representation-Enabled Enforcement (P2)}\label{heading:research-representation-enabled-enforcement}

\emph{Hypothesis H2 (Representation-enabled enforcement). For a declared set of checkable obligations, independent model-derived enforcement reduces the escape of covered violations compared with producer guidance alone, with existing checks held fixed.}

% <!-- point: fixed-gate-comparison-and-factorial-extension | Crossing producer access to the model with independent enforcement separates effects on fulfillment from effects on admission, including their costs and interaction. -->

A four-condition experiment crosses producer access to the model with independent enforcement of its obligations. Hold the task, producer, existing checks, and independent acceptance criteria fixed. Giving the producer the model may change the likelihood of a satisfactory candidate; adding a validator and gate may change which candidates are admitted. Measure correct delivery, covered violations in proposed and admitted work, false rejection, rejected attempts, repair, and human review. This separates improved production from improved admission and exposes the cost of achieving either.

% Evidence: E15 | source: book/data/local-evidence/evidentiary-basis-11202026/integration/records/E15.md | kind: implementation

DocAble's controlled-mutation interface suggests a concrete test: vary whether a reasoner receives the typed-operation model and whether an independent applier enforces selected guards. Evaluate the resulting mutations independently. A reduction in escaped failures could otherwise conceal a mechanism that rejects most work, while apparent compliance with the guards could conceal obligations they do not cover.

For a sound gate over a covered property, exclusion follows from its assumptions. Trials can expose violations of those assumptions and measure practical coverage and cost; they cannot by themselves prove universal soundness. H2 would be weakened if the added enforcement does not reduce covered escape, while false rejection or disproportionate repair could make the intervention economically unattractive even when it does. Semantic obligations outside the check's scope remain part of the independent assessment.

\section{Testing Governance Conversion (P3)}\label{heading:governance-and-engineering-capital}

\emph{Hypothesis H3 (Governance conversion). For recurring failure classes, preserving the relevant engineering judgment in reusable models, procedures, instructions, or controls reduces recurrence or repeated reconstruction and intervention in subsequent comparable work relative to local repair alone.}

% <!-- point: event-centered-design-tests-conversion-and-propagation | Following a failure class can reveal inheritance, but comparisons with local repair or later work without the originator are needed to distinguish conversion from learning or changed tasks. -->

An event-centered longitudinal study follows a failure class from first occurrence through diagnosis, local repair, any durable conversion, later reuse, and maintenance or retirement. Record what the intervention enables as well as which failures recur: a model may make larger changes tractable, while a procedure may let another engineer handle an exception. Instructions qualify when their later use and effects can be observed; conversion need not produce an executable gate.

A decline in recurrence alone cannot establish its cause. The engineer may have become more experienced, or the work may have changed. Compare similar work receiving local repair or on-demand reconstruction, and observe whether benefits persist without the originating engineer. Continued benefit under new personnel or agents is evidence that the environment preserves useful judgment. Repeated dependence on the originator would weaken that interpretation. Across organizations, the same design can examine how authority connects local diagnosis to changes in the shared environment.

% <!-- point: capital-measurement-includes-upkeep-retirement-and-attribution | Capital comparisons require full cost and continuing fit, and must account for shared or interacting assets instead of assigning every gain to the newest control. -->

Measure construction, carrying cost, and reconciliation or retirement as assets lose fit. Several assets may contribute jointly: a model, typed operation, and admission check can reinforce one another. Estimate their return as a bundle or allocate shared costs and gains explicitly. H3 concerns subsequent capability; whether that gain justifies its cost is a further comparison. An asset may reduce recurrence yet cost more than local repair, making continued investment unwarranted.

\section{Testing Governable Acceleration (P4)}\label{heading:capacity-and-environment}

\emph{Hypothesis H4 (Governable acceleration). For tasks affected by declared architectural obligations, increasing agentic capacity produces larger gains in durable throughput with a validated dependency model and independently enforced admission checks than with source files and ordinary review alone.}

% <!-- point: crossed-capability-and-environment-design-tests-substitution | Cross agentic capacity with independently characterized GEE conditions over matched tasks to distinguish substitution, complementarity and negligible effects. -->

A factorial study crosses foundation-model capability or the number of concurrent agents with specified environmental changes. Vary these forms of capacity separately so their effects remain distinguishable. Unlike P3's study of inheritance over time, this comparison asks how the same capacity increase performs under different arrangements. Measure correct admitted work, escape, intervention, and lifecycle cost over a common horizon.

% <!-- point: independently-assessed-gee-conditions-ground-comparison | A dependency model and admission check can be assessed for correspondence and coverage before both conditions receive the same independent outcome evaluation. -->

H4 names one concrete arrangement. The environments use the same codebase, task population, and production tools; assess the dependency model's correspondence and the check's coverage before measuring outcomes. Both conditions receive the same independent outcome assessment. Equivalent gains from added capacity, or gains favoring ordinary review, would challenge H4 under these conditions. Other task families require other specified investments. These comparisons can reveal substitution between environment and model capability, complementarity, or little environmental contribution without constructing an aggregate GEE score.

% <!-- point: environment-investment-may-dominate-marginal-reasoner-gains | The author conjectures that better environments may offer greater marginal value than stronger reasoners for some task classes; this crossover is not established by controlled evidence. -->

I expect the marginal value of a stronger reasoner eventually to be dominated by the value of a better engineering environment. Indeed, I suspect we have already reached that point for some classes of software-engineering work. The experiments would test this conjecture by comparing investment in more capable agents with investment in the conditions under which they work, including the costs of each arrangement. Repeating the comparison as tasks and GEEs evolve would show whether any advantage persists.

% <!-- point: human-review-compares-fixed-attention-and-misleading-evidence | Compare review surfaces at equal time budgets with access to code, including false reassurance from incorrect obligations or stale evidence. -->

Human review deserves a direct comparison within this design. Give engineers the same review-time budget, but vary whether they begin with code or with changed models and verification evidence, retaining access to code when needed. Include cases where a passing check rests on an incorrect obligation or stale evidence; a reassuring display can make the wrong decision easier. As agent output increases, measure whether engineers still distinguish acceptable changes from faulty ones and how much correct work is admitted. Faster approval alone would not support P4 if more violations enter the system.

\section{Engineering Routing}\label{heading:research-optimizing-the-engineering-environment}

% <!-- point: engineering-routing-chooses-the-production-arrangement | Engineering routing chooses representations, process and assurance as well as the reasoner; a policy can adapt those choices and invest in the inherited GEE. -->

The preceding experiments could guide the selection of an engineering arrangement for a body of work. The economical choice depends on the obligations, the environment already available, the consequences of failure, and the work expected to follow. An investment may benefit later tasks, so the cheapest arrangement for one task need not be the cheapest over a sequence.

\textbf{Engineering routing} selects how a task should be engineered for realization, including its representation, foundation model, tools, iteration, and admission mechanisms. A routing decision may favor a stronger foundation model, a cheaper foundation model with better structure, or a deterministic operation for a sufficiently specified task.

% <!-- point: engineering-investment-is-a-constrained-sequential-decision | An engineering policy minimizes lifecycle cost over declared work and a horizon, subject to delivery, heterogeneous assurance and human-attention requirements as the GEE changes. -->

Let \(\pi\) be an engineering policy that selects a production arrangement and any investment in the GEE for each task. Its choices depend on the inherited environment, the available evidence, and current conditions, including the capabilities and prices of available foundation models. Let \(\Pi\) be the set of policies available to the engineer. For a declared task sequence and evaluation horizon of \(H\) engineering episodes, let \(C_i\) denote the total cost incurred in episode \(i\): production, checking, repair, valued human attention, and investment or maintenance attributable to that episode. Charge initial construction to the first episode and allocate shared costs once across the horizon. One formulation of the investment problem is

\[
\begin{aligned}
\min_{\pi\in\Pi}\quad& \mathbb{E}_{\pi}\left[\sum_{i=1}^{H} C_i\right]\\
\text{subject to}\quad& \mathbb{E}_{\pi}[D_H]\ge D_{\min},\\
& \mathcal{A}(\pi,H)=\text{acceptable},\\
& \mathbb{E}_{\pi}[A_H]\le A_{\max},\\
& \mathcal{G}_{i+1}=f(\mathcal{G}_i,a_i),\quad i=1,\ldots,H.
\end{aligned}
\]

Here \(D_H\) counts independently assessed correct admitted delivery and \(D_{\min}>0\) is the minimum required delivery. \(A_H\) is total human attention consumed over the horizon, measured in a declared unit such as engineer-hours, and \(A_{\max}\) is the available attention budget. The expectation is over the outcomes under policy \(\pi\). Human attention is both a valued cost and a limited resource; its constraint does not add a second charge. The condition \(\mathcal{A}\) represents the declared assurance obligations, which may combine invariants, probabilistic failure bounds, and required human authorization. These obligations must also specify which requests or task classes require delivery when an aggregate count would permit the policy to select only easy work. The positive delivery minimum rules out rejecting everything; coverage requirements address selective avoidance.

The final relation retains the schematic inheritance model: episode work \(a_i\), including any engineering investment, changes the GEE available to later work. The policy generalizes the earlier fixed-foundation-model comparison by allowing the reasoning engine and production arrangement to change between episodes. This is an illustrative problem formulation; particular settings will require their own assurance conditions, resource budgets, and valuations of useful delivery.

% <!-- point: optimal-investment-is-unknown-and-conditions-move | The policy, costs and GEE transitions are uncalibrated; changes in foundation models, tasks and obligations can change which engineering investments pay. -->

We do not know the optimal policy. Cost and transition functions remain uncalibrated, and the effects of representations and controls depend on the tasks. Changing foundation-model capabilities and prices may favor reconstructing some knowledge on demand while making other representations economical to maintain. Requirements and organizational authority change as well. The experiments can begin estimating these relationships so that routing decisions account for the continuing value of the environment they leave behind.

% <!-- point: research-needs-durable-outcomes-not-activity-proxies | Effects are limited to the studied tasks, agents, and environments; replication and evidence of adequate representation and evaluation are needed for broader transfer. -->

Effects established for particular tasks, agents, and environments require replication before broader transfer can be justified. Within those limits, evidence can guide where to invest in reusable engineering structure, where to rely on probabilistic reasoning, and where human judgment remains necessary. The economically appropriate division of labor can change as the technology and the work change.
